% test-007.tex -- tagged-PDF test de llaves y valores adicionales
%
% Compilar:  lualatex-dev test-007.tex
% Validar:   show-pdf-tags --xml test-007.pdf | rnv-wrapp
%            verapdf --flavour ua2 --format text test-007.pdf
%
% Test de prueba de los valores de llaves y de las llaves por comando
% que no cubren los tests 001 a 006: tamanos de delimitadores, simbolo y
% subindice de la circunferencia, medidas estadisticas, cifras-sep en
% otros comandos y \setspintent con valores validos. Los "doc:" son
% lecturas reales de NVDA 2026.2 con MathCAT.
\DocumentMetadata
  {
    lang=es-CL,
    pdfversion=2.0,
    pdfstandard=ua-2,
    tagging=on,
    tagging-setup={math/setup={mathml-SE}},
  }
\documentclass{article}
\usepackage[spanish]{babel}
\usepackage{unicode-math}
\usepackage{spintent}
\usepackage{hyperref}
\hypersetup
  {
    pdfdisplaydoctitle = true,
    pdftitle           = {test-007: llaves y valores adicionales},
  }

\begin{document}

\section*{A. Tamaño de los delimitadores}

% doc: «intervalo cerrado de 1 a 5»
1. Intervalo sin tamaño: $\spinterval[size=none]{[1,5]}$.

% doc: «intervalo cerrado de 1 a 5»
2. Intervalo con delimitadores \texttt{big}: $\spinterval[size=big]{[1,5]}$.

% doc: «intervalo cerrado de 1 a 5»
3. Intervalo con delimitadores \texttt{bigg}: $\spinterval[size=bigg]{[1,5]}$.

% doc: «intervalo cerrado de 1 a 5»
4. Intervalo con delimitadores \texttt{Bigg}: $\spinterval[size=Bigg]{[1,5]}$.

% doc: «valor absoluto de x»
5. Valor absoluto sin tamaño: $\spabs[size=none]{x}$.

% doc: «valor absoluto de x»
6. Valor absoluto con barras \texttt{big}: $\spabs[size=big]{x}$.

% doc: «valor absoluto de x»
7. Valor absoluto con barras \texttt{bigg}: $\spabs[size=bigg]{x}$.

% doc: «valor absoluto de x»
8. Valor absoluto con barras \texttt{Bigg}: $\spabs[size=Bigg]{x}$.

\section*{B. Circunferencia: símbolo y subíndice}

% doc: «circunferencia de centro en O»
9. Circunferencia con símbolo \texttt{odot}: $\spcirc[print-sym=odot](O)$.

% doc: «área de la circunferencia con centro en O»
10. Área con símbolo \texttt{odot}: $\spAcirc[print-sym=odot](O)$.

% doc: «perímetro de la circunferencia con centro en O»
11. Perímetro con símbolo \texttt{odot}: $\spPcirc[print-sym=odot](O)$.

% doc: «circunferencia sub 1 de centro en O»
12. Circunferencia con subíndice leído: $\spcirc[sub-index={1},read-sub=true](O)$.

% doc: «área de la circunferencia sub 1 con centro en O»
13. Área con subíndice leído: $\spAcirc[sub-index={1},read-sub=true](O)$.

% doc: «perímetro de la circunferencia sub 1 con centro en O»
14. Perímetro con subíndice leído: $\spPcirc[sub-index={1},read-sub=true](O)$.

\section*{C. Llaves de las medidas estadísticas}

% doc: «moda de x»
15. Moda con subíndice: $\spmoda[sub-index={x}]$.

% doc: «moda sub x»
16. Moda con subíndice leído: $\spmoda[sub-index={x},read-sub=true]$.

% doc: «mediana de x»
17. Mediana con subíndice: $\spmediana[sub-index={x}]$.

% doc: «rango de x»
18. Rango con subíndice: $\sprango[sub-index={x}]$.

% doc: «amplitud»
19. Rango con verbalización propia: $\sprango[read-arg={amplitud}]$.

% doc: «(silencio)»
20. Rango intercuartil silenciado: $\spriq[silent=true]$.

% doc: «mínimo muestral»
21. Mínimo con verbalización propia: $\spmin[read-arg={mínimo muestral}]$.

% doc: «(silencio)»
22. Máximo silenciado: $\spmax[silent=true]$.

% doc: «varianza»
23. Varianza con verbalización propia: $\spvarz[read-arg={varianza}]$.

% doc: «varianza s de x al cuadrado»
24. Varianza con subíndice sin leer: $\spvarz[sub-index={x},read-sub=false]$.

% doc: «(silencio)»
25. Desviación estándar silenciada: $\spdesv[silent=true]$.

% doc: «primer decil»
26. Decil con verbalización propia: $\spdecil[read-arg={primer decil}]{1}$.

% doc: «(silencio)»
27. Percentil silenciado: $\spPerc[silent=true]{25}$.

% doc: «percentil veinticinco»
28. Percentil con verbalización propia: $\spPerc[read-arg={percentil veinticinco}]{25}$.

\section*{D. Separación de cifras en otros comandos}

% doc: «1000000 de pesos»
29. Dinero sin separación de cifras: $\spmoney[cifras-sep=false]{1000000}$.

% doc: «1000000 metros»
30. Unidad sin separación de cifras: $\spunit[cifras-sep=false]{1000000 m}$.

% doc: «1000000 dividido 2»
31. División sin separación de cifras: $\spdiv[cifras-sep=false]{1000000}{2}$.

% doc: «1000000 por 2»
32. Multiplicación sin separación de cifras: $\spmul[cifras-sep=false]{1000000}{2}$.

\section*{E. El comando \texttt{\textbackslash setspintent}}

% doc: «media aritmética de x»
33. Media con subíndice configurado: {\setspintent{spmedia}{sub-index={x}}$\spmedia$}.

% doc: «media aritmética»
% en AF: «media aritmética de x» (misma fuente que el caso 33)
34. Media sin configurar, después del grupo: $\spmedia$.

% doc: «23 metros partido por segundos cuadrado»
35. Unidad con barra configurada: {\setspintent{\spunit}{read-slash=partido-por}$\spunit{23 m/s^{2}}$}.

% doc: «23 metros por segundos cuadrado»
% en AF: «23 metros partido por segundos cuadrado» (misma fuente que el 35)
36. Unidad sin configurar, después del grupo: $\spunit{23 m/s^{2}}$.

% doc: «23 negativo»
37. Signo leído configurado: {\setspintent{\spnum}{read-sign=clear}$\spnum{-23}$}.

% doc: «(silencio)»
38. Probabilidad silenciada por configuración: {\setspintent{spProb}{silent=true}$\spProb{A}$}.

% doc: «20 dividido por 5»
39. Lectura del dividido configurada: {\setspintent{\spdsym}{read-sym=dividido-por}$20 \spdsym 5$}.

% doc: «mcd entre 2 3 y 6»
40. Lectura abreviada configurada: {\setspintent{spmcd}{read-sty=terse}$\spmcd(2, 3, 6)$}.

\end{document}
